\section{Limitations}\label{sec:limitations}

\begin{enumerate}
  \item \textbf{Bounded ROI required.}  LEO-DTF assumes the terminal location
        is approximately known to lie within a predefined region of interest.
        If the true position lies outside the grid, the estimator will return
        the nearest grid cell and no alarm is raised.  Adaptive or hierarchical
        grid strategies are left to future work.

  \item \textbf{Not a GNSS replacement.}  The method provides coarse
        localization and uplink state estimation only.  Centimeter-level or
        meter-level accuracy is not claimed and is not achievable with the
        current approach.

  \item \textbf{Not meter-level localization.}  CRLB analysis shows that
        sub-kilometer localization accuracy requires timestamp precision of
        $\lesssim 10$ $\mu$s (corresponding to $\lesssim 3$ km range noise), which
        is challenging for low-cost IoT hardware.  The current work targets
        ROI-reduction rather than precise geolocation.

  \item \textbf{No real satellite OTA validation.}  All evaluations to date
        use synthetic passes.  The trace-driven HIL plan describes a path to
        hardware validation but is not yet executed.  No claims of
        over-the-air performance are made or implied by the current repository
        state.

  \item \textbf{Pass geometry sensitivity.}  Estimation accuracy depends on
        the satellite elevation trajectory.  Short passes at low maximum
        elevation and nearly overhead passes with symmetric geometry provide
        weaker horizontal position information.  The CRLB sensitivity diagnostic
        characterizes this quantitatively.

  \item \textbf{Oscillator quality sensitivity.}  Large CFO and drift increase
        the correlation between nuisance parameters and position, degrading
        the posterior peak.  The estimator's prior regularization on $b_0$ and
        $b_1$ mitigates this to some extent but cannot fully compensate for
        very poor oscillator stability.

  \item \textbf{Delay observation availability.}  Packet arrival timestamps
        may not be reliably extracted from low-cost LR-FHSS or LoRa receivers.
        If delay observations are absent, the estimator relies solely on
        Doppler, increasing the CRLB and reducing accuracy.  Current LoRa and
        LR-FHSS feature extractors are stubs (see
        Sections~\ref{sec:hil} and~\ref{sec:evaluation}).

  \item \textbf{Grid resolution and computation.}  The grid estimator scales as
        $O(G \cdot D \cdot N)$, where $G$ is the number of position grid
        points, $D$ the number of time-offset grid points, and $N$ the number
        of observations.  Finer grids increase accuracy at increased
        computational cost.  Hierarchical or multi-resolution approaches are
        not yet implemented.

  \item \textbf{HIL is trace-driven, not operational satellite validation.}
        The HIL plan describes offline injection of Doppler/CFO/delay into IQ
        captures; it does not involve real satellite links.  Validating the
        full signal-processing pipeline from real captures remains future work.

  \item \textbf{LoRa / LR-FHSS feature extractors are stubs.}
        \texttt{src/leodtf/feature\_extract\_lora.py} and
        \texttt{src/leodtf/feature\_extract\_lrfhss.py} are placeholders.  Real
        extraction algorithms for packet timing and frequency from LoRa and
        LR-FHSS waveforms are left to future HIL implementation.
\end{enumerate}